% Indicate the main file. Must go at the beginning of the file.
% !TEX root = ../main.tex

%----------------------------------------------------------------------------------------
% CHAPTER 1
%----------------------------------------------------------------------------------------

\chapter{Introduction and Theoretical Background} % Main chapter title
\label{Chapter1} % For referencing the chapter elsewhere, use \ref{Chapter1} 

%----------------------------------------------------------------------------------------

% Define some commands to keep the formatting separated from the content
% Placing such commands in the preamble is a good idea.
\newcommand{\keyword}[1]{\textbf{#1}}
\newcommand{\tabhead}[1]{\textbf{#1}}
\newcommand{\code}[1]{\texttt{#1}}
\newcommand{\file}[1]{\texttt{\bfseries#1}}
\newcommand{\option}[1]{\texttt{\itshape#1}}
%----------------------------------------------------------------------------------------
% SECTION
%----------------------------------------------------------------------------------------

\section{Somatic SNVs Definition and Relevance}

\section{Somatic SNVs: Definition and Relevance}

Genetic variation is the molecular foundation of biological diversity, adaptation, and disease. 
At the broadest level, two fundamental classes of genetic variation shape phenotypic diversity and disease susceptibility in any organism: those inherited through the germline, and those acquired somatically during an individual's lifetime.

\keyword{Germline variants} are present from conception, encoded in every cell of the organism, and passed from parent to offspring through reproduction. Their population frequencies are shaped over generations by evolutionary forces such as natural selection, genetic drift, and recombination, and they underlie a broad spectrum of heritable conditions, from monogenic disorders to complex polygenic disease risk \parencite{YuGeneticvariation2024}.

\keyword{Somatic mutations}, by contrast, do not follow Mendelian inheritance. 
They arise during an individual's lifetime in individual somatic cells as a consequence of DNA damage or errors in DNA replication, and are propagated exclusively to the clonal descendants of the founding cell \parencite{YuGeneticvariation2024}. Because they originate in a single cell and spread only through cell division, somatic mutations are present in only a subset of an organism's cells.

Genetic variants can be divided into four principal classes: \keyword{single-nucleotide variants} (SNVs), which are point substitutions at individual genomic positions; short insertions and deletions (\keyword{indels}, $<$50~bp); \keyword{structural variants} (SVs), encompassing large deletions, duplications, inversions, and translocations; and large \keyword{chromosomal abnormalities}, such as whole-chromosome gains or losses \parencite{Xureviewsomatic2018,YuGeneticvariation2024}.
Among these, SNVs are the most abundant and best characterised class. The most common type of SNV, the single-nucleotide polymorphism (SNP), refers to positions at which heritable variation occurs at appreciable frequency in a population; the term SNV is broader, encompassing both germline and somatic single-base substitutions \parencite{ZouSignificanceSingleNucleotide2020}.

This distinction has a direct technical consequence: because a somatic mutation is present in only a fraction of cells, it appears in sequencing data at a low \keyword{variant allele fraction} (VAF). 
This low VAF is both the defining hallmark of somatic mutations and the core technical challenge in theri detection, as sequecing noice and artefacts can occur at comparable frequencies, making confident discrimination between signal and error non-trivial \parencite{Xureviewsomatic2018, MuyasDenovodetection2024}.

%- How do somatic mutations accumulate throughout life and what biological processes drive this?

%- What molecular mechanisms generate somatic SNVs? (replication errors)

%- What is the difference between germline, somatic, and how are they distinguished computationally?

%----------------------------------------------------------------------------------------
% SECTION
%----------------------------------------------------------------------------------------

\section{Bulk and Single-Cell RNA Sequencing Approaches}
- What does bulk RNA-seq measure, and what are its fundamental limitations for cell-level analysis?

- How does scRNA-seq resolve cell-type heterogeneity that bulk RNA-seq cannot?

- What are the two main scRNA-seq paradigms  droplet-based (10X) vs. plate-based (Smart-seq2)  and how do they differ technically?

%----------------------------------------------------------------------------------------
% SECTION
%----------------------------------------------------------------------------------------
\section{Smart-seq2 as a Platform for Variant Calling}
- Plate-based (Smart-seq2) vs. droplet-based (10X)  key differences for variant calling

- Why does full-length transcript coverage in Smart-seq2 create an advantage for variant calling that 3'-end methods cannot provide?

- Why is Smart-seq2 an underexplored platform for variant calling compared to 10X?

%----------------------------------------------------------------------------------------
% SECTION
%----------------------------------------------------------------------------------------
\section{Challenges of SNV Detection from scRNA-seq}
- What is the technical challenge of detecting SNVs from scRNA-seq data?

- What distinguishes single-cell variant calling from bulk variant calling computationally?

- How does the absence of a matched normal sample complicate somatic vs. germline classification?


%----------------------------------------------------------------------------------------
% SECTION
%----------------------------------------------------------------------------------------
\section{Variant Calling Tools}
- What tools were specifically designed for scRNA-seq variant calling (SComatic, Monopogen, cellsnp-lite), and what are their limitations in this context?

- Why did we choose STAR and bcftools? / I Guess it's to early to Answer this questions
%----------------------------------------------------------------------------------------
% SECTION
%----------------------------------------------------------------------------------------
\section{Dataset and Biological Context move to 2.1}
%----------------------------------------------------------------------------------------
% SECTION
%----------------------------------------------------------------------------------------
\section{Research Question and Objectives }
- What is the research question?

- What is the aim of the project?


\subsection{show them!}
